
\documentclass[letterpaper,oneside]{book}%

\newif\ifinstructor
%% uncomment one of the following, depending if you want to include the
%% instructor notes or not
%\instructortrue
\instructorfalse


\newif\ifnotes
%% uncomment one of the following, depending if you want to include the
%% instructor notes or not
%\notestrue
\notesfalse



\usepackage[left=1in,right=2.75in,top=1in,bottom=1in]{geometry}
\marginparwidth 1.75in

\usepackage{rotating}

\usepackage{tabls}
\usepackage{booktabs}
\usepackage{amsmath}
\usepackage{amssymb}
\usepackage{amsthm}
\usepackage{amsfonts}
\usepackage{multicol}
\usepackage{enumitem}
\usepackage{microtype}
\usepackage{tikz}
\usepackage{pgfplots}
\usetikzlibrary{positioning}
\usepackage{multirow}

\usepackage{comment}

\newcommand{\wrapup}{
\bmw{\section{Wrap Up}
Once you have finished the problems in the section and feel comfortable with the ideas, create a short one page lesson plan that contains examples of the key ideas.  You will get a chance to teach from this lesson plan prior to taking the exam. Then log on to I-Learn and download the quiz. Once you have taken the quiz, you can upload your work back to I-Learn and then download the key to see how you did. If you still need to work on mastering some of the ideas, please do so and then demonstrate your mastery though the quiz corrections.}
}
\newcommand{\myscale}{1}

\usepackage{graphicx}
\usepackage{wrapfig}

\newcommand{\ds}{\displaystyle}
\newcommand{\dfdx}[1]{\frac{d#1}{dx}}
\newcommand{\ddx}{\frac{d}{dx}}


\let\oldmarginpar\marginpar
\renewcommand\marginpar[1]{\-\oldmarginpar{\raggedright\footnotesize #1}}
%\renewcommand\marginpar[1]{\-\oldmarginpar[\raggedleft\footnotesize #1]{\raggedright\footnotesize #1}}


%\usepackage[12hr]{datetime}
%\newdateformat{draftdate}{%
%\shortdayofweekname{\THEDAY}{\THEMONTH}{\THEYEAR}, %
%\THEDAY\ \shortmonthname[\THEMONTH] \THEYEAR}
%\draftdate
%\usepackage{eso-pic}
%\AddToShipoutPicture{\put(10,10){\small Draft: \today\ at \currenttime }}%--- version: \MakeUppercase{\svnInfoRevision}}}

\newcommand{\bmw}[1]{#1}
\newcommand{\marginparbmw}[1]{\bmw{\marginpar{#1}}}

% Notes for Larson, 5th edition
\newcommand{\larsonfive}[1]{}

% Instructor-specific material (answers, helps, etc.)
\ifinstructor
  \newcommand{\instructor}[1]{\marginpar{\textbf{Instructor: }#1}}
\else
  \newcommand{\instructor}[1]{}
\fi

\ifnotes
\renewcommand{\thefootnote}{\roman{footnote}}
\newcommand{\note}[1]{\footnote{#1}\marginpar{\fbox{\textbf{\thefootnote}}}}
\else
\newcommand{\note}[1]{}
\fi

\theoremstyle{plain}
\newtheorem{theorem}{Theorem}[chapter]
\newtheorem*{theorem*}{Theorem}
\newtheorem{lemma}[theorem]{Lemma}
\newtheorem*{lemma*}{Lemma}
\newtheorem{proposition}[theorem]{Proposition}
\newtheorem{corollary}[theorem]{Corollary}


\newtheoremstyle{box}%
{}{}% standard spacing before and after
{}% Body style
{}{\bfseries}{.}% Heading indent, font, and punctuation
{ }% space after heading
{\thmname{#1}\thmnumber{ #2}\thmnote{: #3}}% head spec

\newtheoremstyle{problem}%
{}{}% standard spacing before and after
{}% Body style
{}{\bfseries}{}% Heading indent, font, and punctuation
{1em}% space after heading
{\fbox{\thmname{#1}\thmnumber{ #2}\thmnote{: #3}}}% head spec

\theoremstyle{box}
\newtheorem{definition}[theorem]{Definition}
\newtheorem{dfn}[theorem]{Definition}
\newtheorem*{definition*}{Definition}
\newtheorem{observation}[theorem]{Observation}
\newtheorem{remark}[theorem]{Remark}
\newtheorem{example}[theorem]{Example}
\newtheorem{question}[theorem]{Question}
\newtheorem*{prep-problems}{Preparation Problems}

%\newtheorem{problem}[theorem]{Problem}
\theoremstyle{problem}
\newtheorem{problemnum}{Problem}[chapter]
\newtheorem*{problemnum*}{Problem}
\newenvironment{problem}[1][]{\begin{problemnum}[#1]}{\end{problemnum}\nopagebreak\hrule\bigskip}
\newenvironment{problem*}[1][]{\begin{problemnum*}[#1]}{\end{problemnum*}\nopagebreak\hrule\bigskip}
\newtheorem*{reviewnum*}{Review}
\newenvironment{review*}[1][]{\begin{reviewnum*}[#1]}{\end{reviewnum*}\nopagebreak\hrule\bigskip}
\newtheorem*{exerciseenum*}{Exercise}
\newenvironment{exercise*}[1][]{\begin{exerciseenum*}[#1]}{\end{exerciseenum*}\nopagebreak\hrule\bigskip}


% Abbreviations
\newcommand{\ii}{\ensuremath{\vec \imath}}
\newcommand{\jj}{\ensuremath{\vec \jmath}}
\newcommand{\kk}{\ensuremath{\vec k}}
\newcommand{\vv}{\ensuremath{\mathbf{v}}}
\newcommand{\colvec}[1]{\ensuremath{\begin{bmatrix}#1\end{bmatrix}}}
\DeclareMathOperator{\rank}{rank}
\DeclareMathOperator{\rref}{rref}
\DeclareMathOperator{\vspan}{span}
\DeclareMathOperator{\trace}{tr}
\DeclareMathOperator{\proj}{proj}
\DeclareMathOperator{\curl}{curl}
\newcommand{\RR}{\ensuremath{\mathbb{R}}}
% \vp is "vector prime" and corrects spacing when doing something like
% $\vec r'$ (which has the vector and prime almost touching).
% Instead, do something like $\vec r\vp$
\newcommand{\vp}{\ensuremath{^{\,\prime}}}

%The purpose of this code is to allow me to put lines in matrices so that I can create augmented matrices.
\makeatletter
\renewcommand*\env@matrix[1][*\c@MaxMatrixCols c]{%
  \hskip -\arraycolsep
  \let\@ifnextchar\new@ifnextchar
  \array{#1}}
\makeatother

\newcommand{\cl}[1]{  \begin{matrix}  #1  \end{matrix}  }
\newcommand{\bm}[1]{  \begin{bmatrix}  #1  \end{bmatrix}  }
\newcommand{\inv}{^{-1}}
\newcommand{\im}{\ensuremath{\text{im }}}
\newcommand{\R}{\mathbb{R}}
\newcommand{\blank}[1]{\raisebox{0pt}[14pt]{\rule{#1}{1pt}}}

%------------------------------------------------------------------------------------------------------------

\usepackage{color}
\definecolor{darkblue}{rgb}{0, 0, .6}
\usepackage[breaklinks]{hyperref}
\hypersetup{
	colorlinks=true,
	linkcolor=darkblue,
	anchorcolor=darkblue,
	citecolor=darkblue,
	pagecolor=darkblue,
	urlcolor=darkblue,
}

\begin{document}
Integration Unit -- FLOW
\begin{itemize}
\item Prerequisites (Covered in the Course before this Unit)
\begin{itemize}
\item Derivatives (difference quotient, linearization, ...)
\item Discrete Random Variables (pmf function and calculating probabilities with pmf -- also calculation expected value of pmf could be covered be covered before but does not have to be covered before this unit)
\item Student knows what an antiderivative is -- purposely AVOID talking about indefinite integrals until after this unit.
\end{itemize}
\item Big Ideas
\begin{itemize}
\item Day 1 -- Cumulative Distribution Function (CDF) -- for Discrete Random Variables
\item Day 2 -- The ``Rectangle Map" that is the fundamental theorem of calculus (focus on familiar geometric shapes) 
\begin{itemize}
\item Rug function, $R(x)$, divided by area of rug gives Normalized Rug function, $N(x)$.
\item Using geometry and $R(x)$ we can find the CDF, $F(x)$.
\item Define the PDF, $f(x) = F'(x)$.
\item Notice that $F'(x)$ is the same as $N(x)$ (or in other words $f(x) = N(x)$).
\end{itemize}
\item Day 3 -- Finding Inconvenient Areas (Riemann Sums leading to the definition of the definite integral as an area).
\item Day 4 -- Back to the ``Rectangle Map" (even with inconvenient shapes leading to the FTC and connections between derivatives and antiderivatives, antiderivatives and indefinite integrals, and calculating definite integrals using indefinite integrals).
\end{itemize}
\item What's Next?
\begin{itemize}
\item Practice calculating indefinite integrals as antiderivatives.
\item Practice calculating indefinite integrals with CAS.
\item Practice calculating definite integrals with FTC (and CAS as needed).
\item Calculate Expected Value and Variance of Continuous Random Variables.
\end{itemize}
\end{itemize}


\chapter{Probability and Integration}
\section{Day 1}
\begin{problem}
 Suppose we toss a fair coin once, which means the probability that heads or tails appears is equally likely. 
 We will record the number of heads that appears, and let $X$ represent the outcome. 
 \begin{enumerate}
  \item What are the possible outcomes?  $X=0$ or $X=1$. 
  \item What is the probability that $X=0$?   
  \item What is the probability that $X=1$?   
  \item What is the probability that $X=2$?   
  \item We often summarize probability information into a table, when the number of outcomes is finite. 
    $$\begin{array}{|c|c|c|}\hline
    x & 0&1\\\hline
    P(X=x) & \frac{1}{2} & \frac{1}{2}\\\hline
    \end{array}$$
    \item Recall we call this function a probability mass function (pmf) and we can write it as a piece-wise function.
      $$f(x) = 
    \begin{cases}
    \frac{1}{2} & x=0 \\
    \frac{1}{2} & x=1
    \end{cases}.$$
  \item We often graph this information as well, using a point plot. \textcolor{red}{Include a plot here.} 
  \item What is the domain of the pmf?  What is the range of the pmf?
  \item The pmf above shows the probability of single event occurring. Sometimes we want to know the probability that a given value, or less, will occur. This allows us to use any real number as an input. Compute each of the following: 
  \begin{itemize}
   \item $P(X\leq -1)$
   \item $P(X\leq 0)$
   \item $P(X\leq 0.5)$
   \item $P(X\leq 1)$
   \item $P(X\leq 2)$
  \end{itemize}
  We can summarize the information above into a single piecewise defined function, such as
    $$F(x) = 
    \begin{cases}
    0 & x<0 \\
    \frac{1}{2} &0\leq x<1 \\
    1 & x \geq 1
    \end{cases}.$$
    We call this function the cummulative distribution function (cdf) of the random variable $X$. 
  \item  Because the cdf of $X$, $F(x)$, has a domain of all real numbers, the graph is not a point plot. Draw the cdf of $X$ by hand then verify your graph by creating a plot of the function in R.
  \item What is the domain of the cdf?  What is the range of the cdf?  
 \end{enumerate}

\end{problem}

\begin{problem}
 Suppose we toss a fair six-sided die once, which means each number has the same probability of occurring. 
 We will record the value that appears on the die, and let $X$ represent the outcome. 
 \begin{enumerate}
  \item What are the possible outcomes?  
  \item What is the probability that $X=1$?   
  \item What is the probability that $X=2$?   
  \item What is the probability that $X=9$?   
  \item We often summarize probability information into a table, when the number of outcomes is finite. 
    $$\begin{array}{|c|c|c|c|c|c|c|}\hline
    x & 1&2&3&4&5&6\\\hline
    P(X=x) & \frac{1}{6} &\frac{1}{6} &\frac{1}{6} &\frac{1}{6} &\frac{1}{6} & \frac{1}{6}\\\hline
    \end{array}$$
  \item We can also write this pmf as a piece-wise function.
  $$f(x) = 
    \begin{cases}
    \frac{1}{6} & x=1 \\
    \frac{1}{6} & x=2 \\
    \frac{1}{6} & x=3 \\
    \frac{1}{6} & x=4 \\
    \frac{1}{6} & x=5 \\
    \frac{1}{6} & x=6
    \end{cases}.$$
  \item Draw the pmf of $X$, $f(x)$, by hand (then verify your graph by creating a plot of the function in R).
  \item What is the domain of the pmf?  What is the range of the pmf?
  \item The pmf above shows the probability of single event occurring. Sometimes we want to know the probability that a given value, or less, will occur. This allows us to use any real number as an input. Compute each of the following: 
  \begin{itemize}
   \item $P(X\leq 0)$
   \item $P(X\leq 2)$
   \item $P(X\leq 4.5)$
   \item $P(X\leq 5.1)$
   \item $P(X\leq 8)$
  \end{itemize}
  We can summarize the information above into a single piecewise defined function, such as
    $$F(x) = 
    \begin{cases}
    0 & x<1 \\
    \frac{1}{6} &1\leq x<2 \\
    \frac{2}{6} &2\leq x<3 \\
    \frac{3}{6} &3\leq x<4 \\
    \frac{4}{6} &4\leq x<5 \\
    \frac{5}{6} &5\leq x<6 \\
    1 & x \geq 6
    \end{cases}.$$
    We call this function the cumulative distribution function (cdf) of the random variable $X$. 
  \item  Because the cdf of $X$, $F(x)$, has a domain of all real numbers, the graph is not a point plot. Draw the cdf of $X$ by hand then verify your graph by creating a plot of the function in R.
  \item What is the domain of the cdf?  What is the range of the cdf?  
 \end{enumerate}
\end{problem}



\section{Day 2 -- Prep}
\textcolor{red}{The following problems were inspired by ``Interactive Mathematics Program". Include a reference.}
\begin{problem}
 Let's suppose for a minute that we have a rug (as shown below) and a trap door on the ceiling above the rug. 
 Each time we open the trap door, a dart falls from above to a random spot on the rug. 
 Each point on the rug has the same probability of being hit as any other. 
 When a dart falls on the point $(x,y)$,  we'll record just the $x$-coordinate (so a number between 0 and 5), and let $X$ represent this random variable. 
 
 \begin{tikzpicture}
  \draw (0,0) --node[below]{5 ft} (5,0) -- (5,2) -- (0,2) --node[left]{2ft} (0,0);
 \end{tikzpicture}
 
 \begin{enumerate}

  \item Give 10 different possible outcomes for $X$. How many possible outcomes are there?
  \item Compute each of the following probabilities, and explain how you obtained your answer.
   \begin{enumerate}
    \item $P(0\leq x\leq 2)$
    \item $P(2\leq x\leq 3.5)$
    \item $P(1.8\leq x\leq 5)$
    \item $P(2.263\leq x\leq 3.71)$
   \end{enumerate}
  \item Compute each of the following probabilities. What pattern do you notice? 
   \begin{enumerate}
    \item $P(2\leq x\leq 2.5)$
    \item $P(2\leq x\leq 2.1)$
    \item $P(2\leq x\leq 2.01)$
    \item $P(2\leq x\leq 2.001)$
    \item $P(2\leq x\leq 2.0001)$
   \end{enumerate}
  \item Based on your computations above, what number would you assign to $P(X=2)$?  Explain.  
  \item Is it possible that the dart could fall on a spot whose $x$-coordinate is exactly 2? How likely is this to occur?
   \item Is it possible that the dart could fall on a spot whose $x$-coordinate is exactly 3.224? How likely is this to occur?
  \item Explain why $P(X\leq 2)$ and $P(X<2)$ are equal (also why $P(X\leq 3.224)$ and $P(X<3.224)$ are equal).
  \item Let's obtain the cumulative distribution function. Compute each of the following probabilities. 
   \begin{enumerate}
    \item $P(X\leq 0)$
    \item $P(X\leq 1)$
    \item $P(X\leq 2)$
    \item $P(X\leq 3)$
    \item $P(X\leq 4)$
    \item $P(X\leq 7)$
    \item $P(2\leq X\leq 5)$
    \item $P(2\leq X\leq 6)$
   \end{enumerate}
  \item Draw the cdf for $X$.
  \item Show that for $0\leq x\leq 5$, we have $F(x) = \frac{1}{5}x$.
  \item Show that $$F(x) =\begin{cases}0 & x < 0 \\\frac{1}{5}x & 0 \leq x < 5 \\ 1 & x \geq 5.\end{cases}$$
  \item What is the domain of the cdf?  What is the range of the cdf? 
  \item Explain why $P(a\leq X\leq b) = P(X\leq b)-P(X\leq a)$ (assuming $a$ and $b$ are positive numbers less than 5 with $a<b$).
  \item Explain why $P(a \leq X \leq b) = F(b) - F(a)$.
  \item Show that $P(a \leq X \leq a +\Delta x) = \frac{1}{5} \Delta x$. 
  
  Since $P(a \leq X \leq a +\Delta x) = F(a + \Delta x) - F(a)$ we can rewrite this equation as $F(a + \Delta x) - F(a) = \frac{1}{5} \Delta x$. Notice this last equation is a linearization of the cdf, $F(x)$, at $x = a$.
  \item We can rewrite $F(a + \Delta x) - F(a) = \frac{1}{5} \Delta x$ as $\frac{F(a + \Delta x) - F(a)}{\Delta x} = \frac{1}{5}$. Does anything about this last equation look familiar?
  
  Recall, $\lim_{\Delta x \rightarrow 0}\frac{F(a + \Delta x) - F(a)}{\Delta x} = F'(x)$.
  \item Calculate $F'(x)$. Did you get $F'(x) = \frac{1}{5}$?
 \end{enumerate}
\end{problem}

\begin{problem}
 Imagine now that we have a triangular rug, with a matching triangular trap door above it. Each time we open the trap door, a dart falls from the door onto the rug, and we let $X$ be the random variable that records the $x$-coordinate of the dart's location. 
 \begin{center}
 \begin{tikzpicture}
  \draw (0,0) --node[below]{5 ft} (5,0) --node[right]{5ft} (5,5) -- (0,0);
 \end{tikzpicture}
 \end{center}
 
 \begin{enumerate}

  \item Is $X$ a discrete random variable, or a continuous random variable? Explain?
  \item Compute each of the following probabilities, and explain how you obtained your answer.
   \begin{enumerate}
    \item $P(X\leq 1)$
    \item $P(X\leq 2)$
    \item $P(1\leq X\leq 2)$
    \item $P(X\leq 3)$
    \item $P(2\leq X\leq 3)$
    \item $P(1\leq X\leq 3)$
    \item $P(X \leq 6)$
    \item $P(1.2 \leq 3.4)$
   \end{enumerate}
  \item Explain why $P(a\leq X\leq b) = P(X\leq b)-P(X\leq a)$ (assuming $0 < a < b < 5$, in other words assume $a$ and $b$ are positive numbers less than 5 with $a<b$). 
  \item What is the area of the rug whose $x$-values satisfy $x\leq a$ (assume $a$ is a positive number less than 5).  
  \item What is the area of the entire rug?
  \item Show that $P(0\leq X\leq a) = \frac{a^2}{25}$. 
  \item Why is the CDF of $X$ the function $F(x) = \frac{x^2}{25}$ for $0\leq x\leq 5$. 
  \item Compute $F(5)$.  What does this value mean in the context of this rug problem?
  \item Draw the CDF for $X$.
  \item Show that $$F(x) =\begin{cases}0 &-\infty<x<0\\\frac{1}{25}x^2 & 0\leq x<5 \\ 1 &5\leq x<\infty\end{cases}.$$
  \item Explain why $P(a\leq X\leq b) = F(b)-F(a)$. 
    \item Show that $P(a \leq X \leq a +\Delta x) = \frac{1}{25} (2a + \Delta x) \Delta x$. 
  
  Since $P(a \leq X \leq a +\Delta x) = F(a + \Delta x) - F(a)$ we can rewrite this equation as $F(a + \Delta x) - F(a) = \frac{1}{25} (2a + \Delta x) \Delta x$. Notice this last equation is almost a linearization of the cdf, $F(x)$, at $x = a$.
  \item We can rewrite $F(a + \Delta x) - F(a) = \frac{1}{25} (2a + \Delta x) \Delta x$ as $\frac{F(a + \Delta x) - F(a)}{\Delta x} = \frac{1}{25} (2a + \Delta x)$. Does anything about this last equation look familiar?
  
  Recall, $\lim_{\Delta x \rightarrow 0}\frac{F(a + \Delta x) - F(a)}{\Delta x} = F'(x)$.
  \item Calculate $F'(x)$. Did you get $F'(x) = \frac{1}{25}2a = \frac{2a}{25} = \frac{2}{25} a$?
 
  \item What happened to the $\Delta x$ in the factor $(2a + \Delta x)$?
  \item Notice the left side of the equation $\frac{F(a + \Delta x) - F(a)}{\Delta x}$ has units probability per length. Probability per length is a density, a probability density. So the function $F'(x)$ is usually called a probability density function, $f(x)$. 
  \textcolor{red}{Include a definition of pdf here.}
  \item Notice the function defining the top of the rug is       
  $$R(x) = 
    \begin{cases}
    x & 0 < x < 5 \\
    0 & \text{otherwise}
    \end{cases}.$$
  Notice the total area of the rug (computed in question 5) is $\frac{2}{25}$. How does the function $N(x) = \frac{R(x)}{\text{total rug area}}$ compare to the probability density function $f(x) = F'(x)$?
    
 \end{enumerate}
%I'm thinking of completely changing this section... I think it will be better to show that $P(a\leq x\leq b) = (something)\Delta x$. The width can be factored out, and the term that is left is what we want to talk about.  This motivates introducing the term ``probability density'' naturally, rather than just introducing it artificially. The students can then discover that dividing by the width simplfies things. It also makes the previous problem better. I am liking this. Whe then want to know what happens to the left over terms as we let $\Delta x$ shrink to zero (which we can use software to do if needed, or just do ourselves if it is easy (like a polynomial). ) I don't see any reason to spend gobs of time on limits formally, or even symbolically. We'll be letting things approach zero, and for us in just about all cases that will involve just dropping them off.  Maybe we'll look briefly at the difference quotient for an exponential, and then show how we need a limit of something non-trivial.  But maybe we'll just use software at that point.....  
%In the above computation we found the density per length from $x$ to $x+\Delta x$ to be ....  (Many people call this ratio a difference quotient). Notice that after simplifying the difference quotient, there are two terms. The first term does not depend on the width $\Delta x$, whereas the other part, if we let $\Delta x$ be a really small number, will be extremely small. The first part of the function $\frac{2x}{25}$ has a graph that is a straight line.  If this line represented the top of a triangluar rug, then notice that the rug would have area $A = \frac{1}{2}bh = \frac{1}{2}(5)(\frac{2(5)}{25}) = 1$. The difference quotient used to find probability per area, if we ignore the term that has the $\Delta x$ in it, returns to us the same ideal rug we began with. (I need a problem before this that introduces an ``ideal rug'', namely a rug that has area 1.)    
\end{problem}

\section{Day 2}
Let $R(x)$ be a function that defines the top of a rug and each rug have a matching triangular trap door above it. Each time we open the trap door, a dart falls from the door onto the rug, and we let $X$ be the random variable that records the $x$-coordinate of the dart's location. For each of the rug function, $R(x)$, answer the following questions.

\begin{enumerate}
\item Find $P(0 \leq X \leq a)$, assume $0 < a < 5$.
\item Find the CDF, $F(x)$.
\item Graph $F(x)$.
\item Show that $P(a \leq X \leq a + \Delta x) = \text{Fill in the blank.} \Delta x$.
\item Since $P(a \leq X \leq a + \Delta x) = F(a + \Delta x) - F(a)$, rewrite the above equation and to find $F'(x) = \frac{F(a + \Delta x) - F(a)}{\Delta} = \text{Fill in the blank.}$.
\item Calculate the derivative of the CDF, $F(x)$. Verify you get the same function for $F'(x)$ with both calculations.
\item What is the total area of the rug?
\item Divide the rug function, $R(x)$, by the total area of the rug to get $N(x)$. What do you notice about $F'(x)$ and $N(x)$?
\end{enumerate}

Rug Functions:
$$R_1(x) = 
    \begin{cases}
    -x + 5 & 0 < x < 5 \\
    0 & \text{otherwise}
    \end{cases}.$$

$$R_2(x) = 
    \begin{cases}
    2 & 0 < x < 3 \\
    5 - x & 3 \leq x < 5 \\
    0 & \text{otherwise}
    \end{cases}.$$

$$R_3(x) = 
    \begin{cases}
     & 0 < x < \frac{5}{2} \\
     & \frac{5}{2} \leq x < 5 \\
    0 & \text{otherwise}
    \end{cases}.$$

$$R_4(x) = 
    \begin{cases}
    5 & 0 < x \leq 1 \\
    4 & 1 < x \leq 2 \\
    3 & 2 < x \leq 3 \\
    2 & 3 < x \leq 4 \\
    1 & 4 < x \leq 5 \\
    0 & \text{otherwise}
    \end{cases}.$$

$$R_5(x) = 
    \begin{cases}
    1 & 0 < x \leq 1 \\
    2 & 1 < x \leq 3 \\
    3 & 3 < x \leq 5 \\
    0 & \text{otherwise}
    \end{cases}.$$

$$R_6(x) = 
    \begin{cases}
    25 & 0 < x \leq 1 \\
    24 & 1 < x \leq 2 \\
    21 & 2 < x \leq 3 \\
    16 & 3 < x \leq 4 \\
    9 & 4 < x \leq 5 \\
    0 & \text{otherwise}
    \end{cases}.$$

$$R_7(x) = 
    \begin{cases}
    25 & 0 < x \leq \frac{1}{2} \\
    \frac{99}{4} & \frac{1}{2} < x \leq 1 \\
    24 & 1 < x \leq \frac{3}{2} \\
    \frac{91}{4} & \frac{3}{2} < x \leq 2 \\
    21 & 2 < x \leq \frac{5}{2} \\
    \frac{75}{4} & \frac{5}{2} < x \leq 3 \\
    16 & 3 < x \leq \frac{7}{2} \\
    \frac{51}{4} & \frac{7}{2} < x \leq 4 \\
    9 & 4 < x \leq \frac{9}{2} \\
    \frac{19}{4} & \frac{9}{2} < x \leq 5 \\
    0 & \text{otherwise}
    \end{cases}.$$
    
$$R_8(x) = 
    \begin{cases}
    25-x^2 & 0 < x < 5 \\
    0 & \text{otherwise}
    \end{cases}.$$



\section{Day 3 -- Prep}
$$R_1(x) = 
    \begin{cases}
    25-x^2 & 0 < x < 5 \\
    0 & \text{otherwise} \\
    \end{cases}.$$

$$R_3(x) = 
    \begin{cases}
    x^2 - 10x + 25 & 0 < x < 5 \\
    0 & \text{otherwise} \\
    \end{cases} = \\
    \begin{cases}
    (x-5)^2 & 0 < x < 5 \\
    0 & \text{otherwise}
    \end{cases}.$$

\section{Day 3}
Riemann Sums



\section{Day 4 -- Prep}

\section{Day 4}
Like Day 2 but with ``non-geometric" (or less familiar shapes)

% \begin{comment}
% \end{comment}
\end{document}
